\section{Random Sampling and the Central Limit Theorem}

\subsection{Motivation: Estimating Population Parameters}

Suppose we try to find the average age in a city, say Oaxaca. One way to do this would be by \emph{brute force}—that is, collecting this information person by person. But this method would be very costly in terms of infrastructure and time.

In statistics, this is a common problem whose solution lies in \emph{random sampling}: We take a group of 1,000 individuals (or 10,000 depending on capacity; obviously, the more, the better) and calculate the average age in this group, which we denote by $A_{1}.$

We repeat this procedure, say 100 times, and denote by $A_{1}, A_{2},...,A_{100}$ the average age obtained in each respective attempt.

\subsection{Law of Large Numbers}

According to the \emph{Law of Large Numbers}, the quantity
\begin{align}
	\bar{A}_{100}=\dfrac{A_{1}+...+A_{100}}{100}
\end{align}
is a very close approximation to the true average age of the inhabitants of the city.

\begin{remark}
	We are no longer interested in obtaining the exact value of the average age, but rather in establishing an \emph{estimator} for it. In such a case, we must settle for defining a \emph{range of values} in which the true value might lie.
\end{remark}

\subsection{Central Limit Theorem}

According to the \emph{Central Limit Theorem}, if the number of such samples is sufficiently large, $A_{1},A_{2},...,A_{100}$ will be normally distributed.

\begin{theorem}[Central Limit Theorem]
	If $X_1, X_2, \ldots, X_n$ are independent and identically distributed random variables with mean $\mu$ and variance $\sigma^2$, then when $n$ is sufficiently large:
	\begin{align}
		\frac{\bar{X} - \mu}{\sigma/\sqrt{n}} \xrightarrow{d} N(0,1)
	\end{align}
	
	That is, the distribution of the standardized sample mean approaches a standard normal distribution.
\end{theorem}

\subsection{Practical Implications of the CLT}

The Central Limit Theorem has profound implications for statistical inference:

\begin{enumerate}
	\item \textbf{Independence from the original distribution:} It does not matter whether the original population follows a normal, exponential, uniform, or any other distribution. The distribution of the sample mean will be approximately normal for large samples.
	
	\item \textbf{Sample size:} In practice, $n \geq 30$ is usually sufficient for the normal approximation to be good, although this depends on the shape of the original distribution.
	
	\item \textbf{Parameters of the sampling distribution:}
	\begin{align}
		E(\bar{X}) &= \mu \\
		\text{Var}(\bar{X}) &= \frac{\sigma^2}{n} \\
		\text{SD}(\bar{X}) &= \frac{\sigma}{\sqrt{n}}
	\end{align}
\end{enumerate}

\subsection{Illustrative Example}

Let us consider rolling a fair die. The distribution of a single roll is discrete uniform on $\{1, 2, 3, 4, 5, 6\}$ with:
\begin{align}
	\mu &= \frac{1+2+3+4+5+6}{6} = 3.5 \\
	\sigma^2 &= \frac{(1-3.5)^2 + (2-3.5)^2 + \cdots + (6-3.5)^2}{6} = \frac{35}{12} \approx 2.917
\end{align}

If we roll the die $n = 36$ times and calculate the average $\bar{X}$, by the CLT:
\begin{align}
	\bar{X} \sim N\left(3.5, \frac{2.917}{36}\right) = N(3.5, 0.081)
\end{align}

This means that:
\begin{align}
	P(3.0 < \bar{X} < 4.0) &= P\left(\frac{3.0-3.5}{\sqrt{0.081}} < Z < \frac{4.0-3.5}{\sqrt{0.081}}\right) \\
	&= P(-1.76 < Z < 1.76) \\
	&\approx 0.921
\end{align}

There is approximately a 92.1\% probability that the average of 36 rolls lies between 3.0 and 4.0.

\subsection{CLT Simulation with Python}

We can verify the Central Limit Theorem through simulation:

\begin{lstlisting}[language=Python]
import numpy as np
import matplotlib.pyplot as plt

# Parameters
n = 30  # sample size
num_simulations = 10000

# Simulation: average of n values from an exponential distribution
sample_means = []
for _ in range(num_simulations):
    sample = np.random.exponential(scale=2.0, size=n)
    sample_means.append(np.mean(sample))

# Plot histogram
plt.hist(sample_means, bins=50, density=True, alpha=0.7)

# Overlay theoretical normal curve
mu = 2.0  # mean of the exponential distribution
sigma = 2.0  # standard deviation of the exponential distribution
x = np.linspace(mu - 4*sigma/np.sqrt(n), mu + 4*sigma/np.sqrt(n), 100)
y = (1/(sigma/np.sqrt(n) * np.sqrt(2*np.pi))) * np.exp(-0.5*((x-mu)/(sigma/np.sqrt(n)))**2)
plt.plot(x, y, 'r-', linewidth=2)

plt.xlabel('Sample Mean')
plt.ylabel('Density')
plt.title(f'Distribution of Sample Means (n={n})')
plt.show()
\end{lstlisting}

\subsection{Applications of the CLT}

The Central Limit Theorem is fundamental in many areas:

\begin{itemize}
	\item \textbf{Quality control:} Means of quality measurements follow normal distributions, allowing the use of control charts.
	
	\item \textbf{Surveys and polling:} Sample proportions (which are averages of binary variables) follow approximate normal distributions.
	
	\item \textbf{Finance:} Average returns on investment portfolios are modeled as normal.
	
	\item \textbf{Natural sciences:} Many biological and physical measurements are the result of many small and independent factors, thus following normal distributions.
\end{itemize}

\subsection{Limitations and Considerations}

Although the CLT is very powerful, it is important to remember:

\begin{enumerate}
	\item \textbf{Sample size:} For highly skewed or heavy-tailed distributions, $n > 30$ may be required for a good approximation.
	
	\item \textbf{Independence:} The observations must be independent. If there is correlation among observations, the CLT may not apply.
	
	\item \textbf{Finite variance:} The CLT requires the population variance to be finite. For distributions with infinite variance (such as the Cauchy distribution), the CLT does not apply.
\end{enumerate}

\subsection{Summary}

Random sampling and the Central Limit Theorem form the backbone of statistical inference:

\begin{itemize}
	\item \textbf{Random sampling} allows us to make inferences about large populations from small samples.
	
	\item The \textbf{Law of Large Numbers} guarantees that our estimates converge to the true values.
	
	\item The \textbf{Central Limit Theorem} allows us to use the normal distribution to make inferences, even when the original population is not normal.
\end{itemize}

These theoretical results justify the estimation techniques and hypothesis tests that we will study in the following sections.
